\section{Slide Highlights}

\subsection{视觉摘录}

Jim Fan 在讲稿中列出了一系列关键幻灯片，以下选取其中的代表页进行解读：

\begin{figure}[p]
\centering
\includegraphics[width=0.95\textwidth,page=2]{slides.pdf}
\caption{Foundation Agent 语义框架与技能库。}
\end{figure}

\begin{figure}[p]
\centering
\includegraphics[width=0.95\textwidth,page=3]{slides.pdf}
\caption{GR00T 的多模态感知管道（Vision + Language + Actuation）。}
\end{figure}

\begin{figure}[p]
\centering
\includegraphics[width=0.95\textwidth,page=5]{slides.pdf}
\caption{Sim-to-Real alignment 的 domain randomization 示意与 guard rail。}
\end{figure}

\begin{figure}[p]
\centering
\includegraphics[width=0.95\textwidth,page=7]{slides.pdf}
\caption{Operational Playbook 与 incident trace 回放流程。}
\end{figure}

\begin{figure}[p]
\centering
\includegraphics[width=0.95\textwidth,page=9]{slides.pdf}
\caption{Governance checklist 与 usage policy 示例。}
\end{figure}

\begin{figure}[p]
\centering
\includegraphics[width=0.95\textwidth,page=11]{slides.pdf}
\caption{Project GR00T 的 multi-robot stack 以及 collab plan。}
\end{figure}

\begin{figure}[p]
\centering
\includegraphics[width=0.95\textwidth,page=12]{slides.pdf}
\caption{LLM + physics solver 融合的 pipeline illustration。}
\end{figure}

\begin{figure}[p]
\centering
\includegraphics[width=0.95\textwidth,page=14]{slides.pdf}
\caption{Replay buffer routing 及 prioritize schema。}
\end{figure}

\begin{figure}[p]
\centering
\includegraphics[width=0.95\textwidth,page=17]{slides.pdf}
\caption{Trace-driven observability board，与 metrics dashboard 结合。}
\end{figure}

\begin{figure}[p]
\centering
\includegraphics[width=0.95\textwidth,page=19]{slides.pdf}
\caption{Pause + analyze 工作流细节与 safety alert icons。}
\end{figure}

\begin{figure}[p]
\centering
\includegraphics[width=0.95\textwidth,page=23]{slides.pdf}
\caption{Future research roadmap：LLM physics、self-supervised、multi-agent。}
\end{figure}

\begin{figure}[p]
\centering
\includegraphics[width=0.95\textwidth,page=25]{slides.pdf}
\caption{Incident response sequence: detect, pause, analyze, recover。}
\end{figure}

\begin{figure}[p]
\centering
\includegraphics[width=0.95\textwidth,page=27]{slides.pdf}
\caption{Reward decomposition chart showing language + physics components。}
\end{figure}

\begin{figure}[p]
\centering
\includegraphics[width=0.95\textwidth,page=30]{slides.pdf}
\caption{Logging taxonomy: sensors, policies, overrides。}
\end{figure}

\begin{figure}[p]
\centering
\includegraphics[width=0.95\textwidth,page=32]{slides.pdf}
\caption{Benchmark ladder comparing dexterity, reliability, generalization。}
\end{figure}

\begin{figure}[p]
\centering
\includegraphics[width=0.95\textwidth,page=34]{slides.pdf}
\caption{Toolchain health dashboard with safety interlocks。}
\end{figure}

\begin{figure}[p]
\centering
\includegraphics[width=0.95\textwidth,page=35]{slides.pdf}
\caption{Data lineage schematic showing robot teleop -> replay -> offline RL。}
\end{figure}

\begin{figure}[p]
\centering
\includegraphics[width=0.95\textwidth,page=36]{slides.pdf}
\caption{Governance maturity curve with audit trail metrics。}
\end{figure}

\begin{figure}[p]
\centering
\includegraphics[width=0.95\textwidth,page=37]{slides.pdf}
\caption{Monitoring dashboard outlining trace, metrics, alerts。}
\end{figure}

\begin{figure}[p]
\centering
\includegraphics[width=0.95\textwidth,page=38]{slides.pdf}
\caption{Incident review template with pause/analyze metrics。}
\end{figure}

\begin{figure}[p]
\centering
\includegraphics[width=0.95\textwidth,page=39]{slides.pdf}
\caption{Lifecycle audit trail showing dataset, config, release pairings。}
\end{figure}

\subsection{本章小结}

幻灯片可视化了 GR00T 的架构、仿真策略、运营 playbook 与治理 checklist，是将文本内容具体化的重要素材。

\newpage

